% ID: SEL-D0024-OA-E02
% ESTADO: validada. S025, cotejada con examen y criterios.
\textbf{a)} Los dos tramos son polinómicos. En $x=2$ el límite
por la izquierda es $0$, y por la derecha es $-4+20-16=0$,
que coincide con $R(2)$. La función es continua en $[1,6]$,
con continuidad lateral en los extremos.

\textbf{b)} El primer tramo va de $(1,-1)$ a $(2,0)$, punto
que incluye el segundo tramo. A partir de ahí sigue un arco
cóncavo con vértice $(5,9)$ y extremo $(6,8)$.
\begin{center}
\setlength{\unitlength}{5mm}
\begin{picture}(12,13)(-1,-2)
\qbezier(0,0)(3.5,0)(7,0)\put(7.2,0){$x$}
\qbezier(0,-1.5)(0,4.5)(0,10.5)\put(0,10.6){$R$}
\qbezier(1,-1)(1.5,-0.5)(2,0)
\qbezier(2,0)(4,12)(6,8)
\put(1,-1.2){\makebox(0,0)[t]{$(1,-1)$}}
\put(2.2,-0.2){$(2,0)$}
\put(5,9.3){\makebox(0,0)[b]{$(5,9)$}}
\put(6.2,8){$(6,8)$}
\end{picture}
\end{center}
\textbf{c)} En el segundo tramo $R(x)=9-(x-5)^2$, de modo
que el máximo se alcanza al invertir \textbf{5 millones de euros}
y asciende a \textbf{9 millones}. El primer tramo es negativo.
En el segundo, $9-(x-5)^2>0$ equivale a $2<x<8$.
Al intersectar con el dominio obtenemos $\boxed{2<x\le6}$.
